% ECONOMICS_PROBLEM: {"id": "G12-481", "title": "Climate/ESG asset pricing and real effects", "jel_code": "G12", "source_id": "OP-0098"}
% !TeX program = xelatex
\documentclass[11pt,a4paper]{article}
\usepackage[margin=23mm]{geometry}
\usepackage{fontspec}
\setmainfont{Latin Modern Roman}
\usepackage{amsmath,amssymb,mathtools}
\usepackage{xcolor,enumitem,booktabs,longtable,array,xurl,hyperref}
\definecolor{muted}{HTML}{53636C}
\hypersetup{colorlinks=true,urlcolor=blue,linkcolor=blue}
\setlength{\parindent}{0pt}
\setlength{\parskip}{5pt}
\newcommand{\R}{\mathbb R}
\newcommand{\E}{\mathbb E}
\newcommand{\N}{\mathbb N}
\newcommand{\ind}{\mathbf 1}
\newcommand{\Var}{\operatorname{Var}}
\newcommand{\Cov}{\operatorname{Cov}}
\newcommand{\supp}{\operatorname{supp}}
\DeclareMathOperator*{\argmin}{arg\,min}
\DeclareMathOperator*{\argmax}{arg\,max}
\newcommand{\entrymeta}[3]{{\sffamily\small\color{muted}\textbf{\texttt{#1}}\quad|\quad #2\par #3\par}}
\newcommand{\Problemref}[1]{\texttt{#1}}
\newcommand{\JELref}[2]{\texttt{#2} (JEL \texttt{#1})}
\begin{document}
\section*{Climate/ESG asset pricing and real effects}
\textbf{Problem ID:} \texttt{G12-481}\quad \textbf{JEL:} \texttt{G12}\par
% BEGIN ECONOMICS STATEMENT
\label{problem:OP-0098}
\label{atlas:prob:F8}

\noindent\textbf{Atlas ID:} F8; \textbf{original number:} 98; \textbf{field:} Finance. \textbf{Classification:} Editorial research formulation (theoretical/empirical); not a certified unsolved theorem.

\subsection*{Definitions and assumptions}
Firms $i$ choose investment $I_i$, abatement $a_i$ and financing within technologies producing cash flows and physical emissions $E_i(I_i,a_i)$. Investors maximize expected utility over budget-feasible portfolios, possibly valuing environmental characteristics directly. Model $\theta$ specifies climate-state risks, tastes, disclosure and rating measurement, ownership and a market-clearing equilibrium selection. For ordinary financial claims, a pricing kernel $m_{t+1}$ satisfies $P_{it}=\mathbb E_t[m_{t+1}X_{i,t+1}]$, with $X_{i,t+1}$ the next-period payoff, when its assumptions apply; investor-specific nonpecuniary benefits or participation constraints require the corresponding modified optimality condition. Assume finite discounted cash flows and expected emissions. Fix a finite horizon $H$ and define $E_{i,\theta}(z)=\sum_{t=0}^{H}E_{it,\theta}(z)$ as total physical emissions in tonnes of carbon-dioxide equivalent under intervention $z$. Investment $I_{i,\theta}(z)$ uses a common monetary numeraire and the same horizon.

\subsection*{Formal research question}
For an externally assigned financing-demand intervention $z\in\{0,1\}$, identify effects on financing costs, investment and emissions:
\[
 \Delta_E=\sum_i\mathbb E_\theta[E_{i,\theta}(1)-E_{i,\theta}(0)],
 \qquad
 \Delta_I=\sum_i\mathbb E_\theta[I_{i,\theta}(1)-I_{i,\theta}(0)].
\]
The summation boundary includes affected competitors or identifies a narrower direct effect explicitly. Decompose observed green price differences into climate-risk compensation, investor tastes and expected cash-flow changes using independently informative data. Determine when shifts in investor demand cause physical mitigation after accounting for new issuance, secondary trading, firms’ substitution and general equilibrium leakage.

\subsection*{Interpretation and scope}
A ``green premium'' needs a sign convention: a higher price usually means a lower expected financial return conditional on comparable payoffs. Ex post green outperformance can arise from unexpected preference or climate news and does not establish higher ex ante required returns. Rating shocks may change information as well as demand, undermining an exclusion restriction based solely on financial costs. Firms can relabel assets, sell emitting operations or shift production, so disclosure improvements are not physical emissions reductions. Costs of capital are endogenous to risk and firm choices; a causal financing channel requires a defensible intervention and mediation assumptions. Welfare requires climate damages and any preference benefits, not the emissions statistic alone. State the intervention's assignment unit and law, consistency between observed and assigned outcomes, and positivity of both assignments on the target state support. A causal design additionally requires random assignment, conditional exchangeability, or another justified identification restriction; market spillovers belong in the assigned regime.

\subsection*{Source audit and status}
[S1] concerns physical climate risks and market efficiency, [S2] carbon emissions and returns, and [S3] equilibrium sustainable investing. They do not jointly identify the causal physical-emissions effect of all ESG finance. This is an editorial asset-pricing and real-effects agenda; no universal sign or currently open theorem is asserted.

\subsection*{References}

\begin{enumerate}[label={[S\arabic*]},leftmargin=*]

\item Harrison Hong, Frank Weikai Li, and Jiangmin Xu (2019). \emph{Climate Risks and Market Efficiency}. Journal of Econometrics 208(1), 265–281. \url{https://www.sciencedirect.com/science/article/pii/S0304407618301817}.
\textit{Locator and role:} Publisher or author abstract / bibliographic record; Physical climate risk, food-sector returns and market efficiency. Provides background or a benchmark result; does not establish that the editorial question remains unresolved in 2026.

\item Patrick Bolton and Marcin Kacperczyk (2021). \emph{Do Investors Care about Carbon Risk?}. Journal of Financial Economics 142(2), 517–549. \url{https://doi.org/10.1016/j.jfineco.2021.05.008}.
\textit{Locator and role:} Publisher or author abstract / bibliographic record; Carbon emissions and stock returns; estimated association does not identify emissions response to financing costs. Provides background or a benchmark result; does not establish that the editorial question remains unresolved in 2026.

\item Ľuboš Pástor, Robert F. Stambaugh, and Lucian A. Taylor (2021). \emph{Sustainable Investing in Equilibrium}. Journal of Financial Economics 142(2), 550–571. \url{https://doi.org/10.1016/j.jfineco.2020.12.011}.
\textit{Locator and role:} Publisher or author abstract / bibliographic record; Equilibrium asset pricing with ESG tastes and climate-risk hedging. Provides background or a benchmark result; does not establish that the editorial question remains unresolved in 2026.

\end{enumerate}

\subsection*{Common conventions: Source mathematical conventions}

Unless an entry states otherwise, agent and firm sets are finite. State and action spaces are measurable, strategies and outcome maps are measurable, probability laws are specified, and expectations exist for the targets under discussion. Infinite-horizon discount factors lie in $(0,1)$ and discounted objectives are finite. Norms, tolerances and complexity input sizes must be specified for computational comparisons. Constants or primitives that appear locally are inputs, not empirical facts asserted by this catalogue.

\textbf{Identification.} A statistical model $m\in\mathcal M$ implies an observable law $P_m$ and target $\theta(m)$. At law $P$, its identified set is
\[
\Theta(P;\mathcal M)=\{\theta(m):m\in\mathcal M,\ P_m=P\}.
\]
Point identification holds when this set is a singleton, or equivalently when $P_{m_1}=P_{m_2}$ implies $\theta(m_1)=\theta(m_2)$ for all compatible models. Sharp bounds mean bounds attained, or approached when the boundary is not attained, by compatible models. Writing this set is a definition, not a solution: the substantive task is to establish informative restrictions and derive or estimate the set. An unrestricted-confounding model may give only trivial bounds.

\textbf{Inference.} Identification, estimability and valid uncertainty quantification are separate questions. Fix $\alpha\in(0,1)$. A confidence procedure $C_n$ over a declared law class $\mathcal P$ has asymptotic uniform coverage at level $1-\alpha$ when
\[
\liminf_{n\to\infty}\inf_{P\in\mathcal P}\Pr_P\{\theta(P)\in C_n\}\geq1-\alpha.
\]
For set-identified targets the coverage event must be defined separately, for example coverage of the full identified set. If an entry uses identification-indexed models instead of uniquely indexed laws, the same coverage convention is applied over those models. Pointwise limits do not establish uniform validity, and asymptotic validity does not guarantee finite-sample precision. Sample sizes, dependence structures and weak-identification sequences are part of each application.

\textbf{Counterfactuals.} $Y_i(a)$ is a potential outcome under a specified intervention, including its timing, implementation and versions. It is not $Y_i$ conditional on voluntarily choosing $a$. A full intervention vector permits interference; shorthand scalar treatment notation does not silently impose no interference. Assignment, selection, attrition, anticipation and measurement must be specified. A claim about an instrument's local effect is not automatically a population policy effect.

\textbf{Economic equilibrium.} A counterfactual model includes best responses, constraints, feasibility, market clearing, state transitions and expectations consistent with the stated information. When several equilibria exist, either a declared selector is an input or the target is set-valued. Comparisons across policy regimes must use the stated selection convention. Reduced-form relationships are not presumed invariant after an equilibrium-changing intervention.

\textbf{Optimization and welfare.} Optimization is over the stated feasible set. If attainment is not established, $\sup$ or $\inf$ and $\varepsilon$-optimal choices replace an asserted maximizer or minimizer. For example, with value $V=\sup_{a\in\mathcal A}W(a)<\infty$, an $\varepsilon$-optimal set is $\{a\in\mathcal A:W(a)\geq V-\varepsilon\}$ for $\varepsilon>0$. Benefits, costs, risk and health or environmental outcomes can be aggregated only using declared units and conversion weights. Interpersonal utility weights, the treatment of fiscal transfers and foreign profits, discounting, and the behavioral welfare criterion are explicit inputs. A comparative-static derivative additionally requires existence and appropriate differentiability of the selected counterfactual.

\textbf{Sources.} Local labels [S1], [S2], and so forth restart in every question. Locators identify the relevant abstract, model, section or result at the level actually checked; a bibliographic or abstract-level verification is not represented as inspection of an entire proof. Working papers are distinguished from later publications and changing versions. Source summaries are paraphrased. All URL checks and editorial formulations refer to the audit date stated above.
% END ECONOMICS STATEMENT
\end{document}
